\section{Related Work}\label{related}

\begin{description}
\item[Proof of Work] Bitcoin was the first successful implementation of a consensus protocol within a decentralized network infrastructure, implementing the novel Proof-of-Work (PoW) protocol, Bitcoin has provided the philosophical foundation for a majority of all subsequent protocols and networks \cite{Nakamoto2008}. The premise of the PoW protocol, separate from the political and economic beliefs embedded within the Bitcoin system, is the belief that a network of pseudonymous and possibly untrustworthy machines can operate in a manner that guarantees the presence of a correct and shared global state. The guarantees of correctness provided by the system are outlined as probabilistic, with a tolerance of \(P<99.99\%\), where $P$ is representative of the probability by which an attacker could “catch up” with the valid network state. This measurement is important as the correct branch of the blockchain is identified by length, with the longest chain perceived as unanimously correct. In contrast to our work on Blossom, as well as the work described in other papers referenced in this section, Bitcoin implements a novel protocol, Proof-of-Work (PoW), which is considered a highly effective yet inefficient protocol. The unique nature of PoW, which has been emulated by subsequent systems, is the competitive nature of the protocol.
\item[Proof of Stake] Unlike the other protocols described here, the cited origin for Proof-of-Stake (PoS) is recognized not as part of a Whitepaper or publication, but as an idea presented by a user within the \textit{bitcointalk forum} \cite{POS}. The protocol, universally seen as the predecessor to Bitcoin’s PoW protocol, is a respected alternative as computation cost and related energy consumption required to validate network transactions in PoS are vastly reduced from PoW. These reductions are immensely important given the important conversations regarding blockchain-related emissions, and additional scalability concerns. The key point of departure in PoS from PoW is the use of randomized leader selection, where the chance of a validators selection as the consensus leader is directly associated with the size of their staking position (staking is generally understood as some deposit of crypto made by a network member). By requiring individuals to deposit a reasonably significant sum of tokens to advance the network, any malicious action can be made accountable by seizing or diluting the staked assets. However, despite the benefits of PoS, there remain additional security and governance issues resulting from the protocol, with the additional fact that the performance of PoS systems remains lower than is necessary for universal blockchain adoption.
\item[Practical Byzantine Fault Tolerance] A highly influential paper describing “a new state-machine replication algorithm,” Practical Byzantine Fault Tolerance (PBFT) is a seminal work that laid the foundation for decentralized consensus protocols \cite{Castro2002}. We frequently reference this work within this paper, as the work done on PBFT was highly influential in the design of Blossom. A characteristic of PBFT we find to be highly valuable is the ability for all network members to participate in consensus. This cooperative behavior provides a deterministic finality to all consensus updates while allowing for all nodes (replica machines) to receive the content of state updates at the time of consensus. While PBFT is designed as a “practical” algorithm, when implemented as the consensus protocol for blockchain networks, the internode communication between members is highly inefficient with a communication overhead of \(\mathcal{O}(n^2)\). When considering the need to handle an exponentially increasing number of communications for every additional node, paired with the increased latency of globally decentralized networks, PBFT fails to be effective at scale. Additional factors which have limited the adoption of PBFT to permissioned networks include the inability of the protocol to successfully mitigate Sybil and DDoS attacks.
\item[HoneyBadgerBFT \& Tendermint] While these two papers describe unique works, both papers propose improvements to the PBFT algorithm, with the goal of optimizing PBFT for use in blockchain systems. Tendermint is generally understood as the most efficient Byzantine fault tolerant (BFT) protocol implementation, with the highest profile implementation of the protocol in the Cosmos network \cite{Buchman2016} \cite{kwon2019cosmos}. We are interested in Tendermint as the proven performance and efficiency improvements displayed in the protocol provide a validated benchmark for us to further improve. HoneyBadgerBFT (HBBFT) is another improvement on the PBFT algorithm; however, unlike Tendermint, HBBFT is designed to increase the attack resistance and asynchrony of the protocol \cite{Miller2016}. Our interest in HBBFT is the result of its protocol architecture, which adds additional tolerances to the protocol, reducing the consequences of singular machine failures. We have made vast improvements beyond both of these algorithms; however, they existed as important reference material in the design of Blossom.
\item[Snowflake] Snowflake is a novel consensus protocol that has been implemented in the Avalanche blockchain \cite{Rocket2018} \cite{allavena2005correctness}. Since the original release of this paper the protocol has seen subsequent improvements; however, for the purpose of our reference, we will be focusing on the original work cited above. To guarantee a shared global state, Snowflake implements a subsampling algorithm with similarities to gossip protocols \cite{allavena2005correctness}, in that each node polls a random subset of the population and reaches a supermajority consensus among said subset. This process is repeated until all nodes in the network return a shared state. The ingenuity of this protocol is the embedded novelty of relying on randomized subsampling, rather than some traditional framework to guarantee a shared network state. However, despite the improvements in communication overhead, Snowflake performance displays constant time performance. It is, for this reason, they heavily cite the implementation of “layer 2 scaling solutions” as the method by which they can further increase network performance.

\item[Proof of History] Another novel protocol, Proof-of-History (PoH) is a highly optimized consensus protocol that utilizes the ordering of transactions in blockchain systems as a proxy for time \cite{yakovenko2018solana}. This protocol is notable for its use in the Solana blockchain, the founders of whom were also the authors of the original PoH paper. While the means by which Solana guarantees ordering remains questionable, the need to order transactions remains one of importance. In previous systems, (i.e Bitcoin) transaction order is a means of validating a shared state, while PoH embeds additional value from this behavior. We reference PoH due to the protocol’s unique behavior, a fact that enticed us into questioning the purpose of transaction ordering in our protocol.
\end{description}